% GENERATED FILE. Edit canonical lessons, then run npm run generate:books.
% canonical-source-hash: fnv1a64:031f63f5fd4c6bf6
% canonical-lessons: RU-C302-zadykhatsya, RU-C302-vyzhivat, RU-C302-perevyazyvat, RU-C302-obzhechsya, RU-C302-porezatsya, RU-R302-review-295-302

\chapter{First aid}
\label{ch:ru-first-aid}
\hlchaptermodality{302}

\begin{chapteropening}
\textbf{By the end of this chapter:} \emph{I can say that someone was choking, survived, had a wound dressed, got burnt or cut himself.}

Complete the last lesson of chapter 302: I can say that someone was choking, survived, had a wound dressed, got burnt or cut himself.
\end{chapteropening}

\section[\texorpdfstring{\ru{задыхаться} (zadykhátsya)}{задыхаться (zadykhátsya)}]{\ru{задыхаться} (zadykhátsya) — to choke, to suffocate; to be out of breath}

\label{lesson:RU-C302-zadykhatsya}



\begin{quote}
Before the new one: say the Russian for to collapse, then the Russian for panic.
\end{quote}

\subsection*{You'll want to know: \ru{задыхаться}}
\textbf{\ru{задыхаться}} (\emph{zadykhátsya}) — \textquotedblleft{}to choke, to suffocate; to be out of breath\textquotedblright{}.

\begin{quote}
\ru{В} \ru{дыму} \ru{мы} \ru{задыхались}.
\end{quote}

\begin{quote}
\emph{V dymú my zadykhális.}
\end{quote}

\begin{quote}
We were choking in the smoke.
\end{quote}

Imperfective; perfective \emph{\ru{задохнуться}} (\emph{zadokhnútsya}). Also after running: \emph{\ru{он} \ru{бежал} \ru{и} \ru{задыхался}}.

\begin{cousinweb}
\emph{\ru{задыхаться}} is built on \emph{\ru{за}-}, \emph{\ru{дыхать}} \textquotedblleft{}to breathe\textquotedblright{} and \emph{-\ru{ся}} (EtymDB), from the Slavic root of \emph{\ru{дышать}} and \emph{\ru{дух}} \textquotedblleft{}breath\textquotedblright{}.
\end{cousinweb}

\begin{grammarlens}[title={What you've built}]
One more word for telling, explaining and describing at B1.
\end{grammarlens}

\subsection*{Your turn}
\begin{itemize}
\raggedright
  \item \emph{Say it:} \emph{zadykhátsya}
  \item \emph{Say it:} \emph{zadykhátsya}, once more
  \item \emph{Say it:} a sentence of your own with \emph{zadykhátsya}
  \item \emph{Recall:} say the Russian for to collapse, then the Russian for panic, then say \emph{zadykhátsya} again
\end{itemize}

\begin{tcolorbox}[breakable,colback=teal!4,colframe=teal!35!black,title={Before you move on}]
What does \ru{задыхаться} mean? (\textquotedblleft{}To choke, to suffocate; to be out of breath\textquotedblright{}.) Say it once more.
\end{tcolorbox}
\section[\texorpdfstring{\ru{выживать} (vyzhivát)}{выживать (vyzhivát)}]{\ru{выживать} (vyzhivát) — to survive}

\label{lesson:RU-C302-vyzhivat}



\begin{quote}
Before the new one: say the Russian for panic, then the Russian for to choke.
\end{quote}

\subsection*{You'll want to know: \ru{выживать}}
\textbf{\ru{выживать}} (\emph{vyzhivát}) — \textquotedblleft{}to survive\textquotedblright{}.

\begin{quote}
\ru{Он} \ru{выжил} \ru{после} \ru{аварии}.
\end{quote}

\begin{quote}
\emph{On výzhil pósle avárii.}
\end{quote}

\begin{quote}
He survived the accident.
\end{quote}

Imperfective; perfective \emph{\ru{выжить}} (\emph{výzhit}), \emph{\ru{выживу}}, stressed on \emph{\ru{вы}-}.

\begin{cousinweb}
\emph{\ru{выживать}} is built on \emph{\ru{вы}-} \textquotedblleft{}out\textquotedblright{} and \emph{\ru{жить}} \textquotedblleft{}to live\textquotedblright{}, an inherited Slavic verb from PIE \emph{g\textsuperscript{w}eyh\textsubscript{3}-} \textquotedblleft{}to live\textquotedblright{}: to live through.
\end{cousinweb}

\begin{grammarlens}[title={What you've built}]
One more word for telling, explaining and describing at B1.
\end{grammarlens}

\subsection*{Your turn}
\begin{itemize}
\raggedright
  \item \emph{Say it:} \emph{vyzhivát}
  \item \emph{Say it:} \emph{vyzhivát}, once more
  \item \emph{Say it:} a sentence of your own with \emph{vyzhivát}
  \item \emph{Recall:} say the Russian for panic, then the Russian for to choke, then say \emph{vyzhivát} again
\end{itemize}

\begin{tcolorbox}[breakable,colback=teal!4,colframe=teal!35!black,title={Before you move on}]
What does \ru{выживать} mean? (\textquotedblleft{}To survive\textquotedblright{}.) Say it once more.
\end{tcolorbox}
\section[\texorpdfstring{\ru{перевязывать} (perevyázyvat)}{перевязывать (perevyázyvat)}]{\ru{перевязывать} (perevyázyvat) — to bandage, to dress (a wound)}

\label{lesson:RU-C302-perevyazyvat}



\begin{quote}
Before the new one: say the Russian for to choke, then the Russian for to survive.
\end{quote}

\subsection*{You'll want to know: \ru{перевязывать}}
\textbf{\ru{перевязывать}} (\emph{perevyázyvat}) — \textquotedblleft{}to bandage, to dress (a wound)\textquotedblright{}.

\begin{quote}
\ru{Медсестра} \ru{перевязала} \ru{рану}.
\end{quote}

\begin{quote}
\emph{Medsestrá perevyazála ránu.}
\end{quote}

\begin{quote}
The nurse dressed the wound.
\end{quote}

Imperfective; the perfective is \emph{\ru{перевязать}} (\emph{perevyazát}), \emph{\ru{перевяжу}}. \emph{\ru{Перевязка}} — a dressing.

\begin{cousinweb}
\emph{\ru{перевязывать}} is built on \emph{\ru{пере}-} \textquotedblleft{}across\textquotedblright{} and \emph{\ru{вязать}} \textquotedblleft{}to tie\textquotedblright{}, inherited from Proto-Slavic \emph{vęzati} (EtymDB): to tie across.
\end{cousinweb}

\begin{grammarlens}[title={What you've built}]
One more word for telling, explaining and describing at B1.
\end{grammarlens}

\subsection*{Your turn}
\begin{itemize}
\raggedright
  \item \emph{Say it:} \emph{perevyázyvat}
  \item \emph{Say it:} \emph{perevyázyvat}, once more
  \item \emph{Say it:} a sentence of your own with \emph{perevyázyvat}
  \item \emph{Recall:} say the Russian for to choke, then the Russian for to survive, then say \emph{perevyázyvat} again
\end{itemize}

\begin{tcolorbox}[breakable,colback=teal!4,colframe=teal!35!black,title={Before you move on}]
What does \ru{перевязывать} mean? (\textquotedblleft{}To bandage, to dress (a wound)\textquotedblright{}.) Say it once more.
\end{tcolorbox}
\section[\texorpdfstring{\ru{обжечься} (obzhéchsya)}{обжечься (obzhéchsya)}]{\ru{обжечься} (obzhéchsya) — to burn oneself, to get burnt}

\label{lesson:RU-C302-obzhechsya}



\begin{quote}
Before the new one: say the Russian for to survive, then the Russian for to bandage.
\end{quote}

\subsection*{You'll want to know: \ru{обжечься}}
\textbf{\ru{обжечься}} (\emph{obzhéchsya}) — \textquotedblleft{}to burn oneself, to get burnt\textquotedblright{}.

\begin{quote}
\ru{Я} \ru{обжёгся} \ru{горячим} \ru{чаем}.
\end{quote}

\begin{quote}
\emph{Ya obzhyógsya goryáchim cháyem.}
\end{quote}

\begin{quote}
I burnt myself on hot tea.
\end{quote}

Perfective; imperfective \emph{\ru{обжигаться}} (\emph{obzhigátsya}). Past \emph{\ru{обжёгся}} (\emph{obzhyógsya}), \emph{\ru{обожглась}}.

\begin{cousinweb}
\emph{\ru{обжечься}} is built on \emph{\ru{об}-} and \emph{\ru{жечь}} \textquotedblleft{}to burn\textquotedblright{}, inherited through Old East Slavic \emph{žeči} (EtymDB), with the reflexive \emph{-\ru{ся}}.
\end{cousinweb}

\begin{grammarlens}[title={What you've built}]
One more word for telling, explaining and describing at B1.
\end{grammarlens}

\subsection*{Your turn}
\begin{itemize}
\raggedright
  \item \emph{Say it:} \emph{obzhéchsya}
  \item \emph{Say it:} \emph{obzhéchsya}, once more
  \item \emph{Say it:} a sentence of your own with \emph{obzhéchsya}
  \item \emph{Recall:} say the Russian for to survive, then the Russian for to bandage, then say \emph{obzhéchsya} again
\end{itemize}

\begin{tcolorbox}[breakable,colback=teal!4,colframe=teal!35!black,title={Before you move on}]
What does \ru{обжечься} mean? (\textquotedblleft{}To burn oneself, to get burnt\textquotedblright{}.) Say it once more.
\end{tcolorbox}
\section[\texorpdfstring{\ru{порезаться} (porézatsya)}{порезаться (porézatsya)}]{\ru{порезаться} (porézatsya) — to cut oneself}

\label{lesson:RU-C302-porezatsya}



\begin{quote}
Before the new one: say the Russian for to bandage, then the Russian for to burn oneself.
\end{quote}

\subsection*{You'll want to know: \ru{порезаться}}
\textbf{\ru{порезаться}} (\emph{porézatsya}) — \textquotedblleft{}to cut oneself\textquotedblright{}.

\begin{quote}
\ru{Он} \ru{порезался} \ru{ножом}.
\end{quote}

\begin{quote}
\emph{On porézalsya nozhóm.}
\end{quote}

\begin{quote}
He cut himself with a knife.
\end{quote}

Perfective. The tool is in the instrumental.

\begin{cousinweb}
\emph{\ru{порезаться}} is built on \emph{\ru{по}-} and \emph{\ru{резать}} \textquotedblleft{}to cut\textquotedblright{}, inherited from Proto-Slavic \emph{rězati} (EtymDB), with the reflexive \emph{-\ru{ся}}.
\end{cousinweb}

\begin{grammarlens}[title={What you've built}]
One more word for telling, explaining and describing at B1.
\end{grammarlens}

\subsection*{Your turn}
\begin{itemize}
\raggedright
  \item \emph{Say it:} \emph{porézatsya}
  \item \emph{Say it:} \emph{porézatsya}, once more
  \item \emph{Say it:} a sentence of your own with \emph{porézatsya}
  \item \emph{Recall:} say the Russian for to bandage, then the Russian for to burn oneself, then say \emph{porézatsya} again
\end{itemize}

\begin{tcolorbox}[breakable,colback=teal!4,colframe=teal!35!black,title={Before you move on}]
What does \ru{порезаться} mean? (\textquotedblleft{}To cut oneself\textquotedblright{}.) Say it once more.
\end{tcolorbox}
\section[(dialogue)]{Words from chapters 295–302 — a review}

\label{lesson:RU-R302-review-295-302}



\begin{quote}
Say the Russian for a fire and smoke.
\end{quote}

\subsection*{Your turn}
Say these aloud:

\begin{itemize}
\raggedright
  \item a fire, smoke, to burn, to catch fire, a fire extinguisher
  \item a rescuer, to save, rescue, to evacuate, an evacuation
  \item an injury, a bruise, consciousness, to be hurt, to wound
  \item a hit, an injured person, to crash, to come out unhurt, to overturn
  \item a stretcher, a siren, an alarm, the person on duty, to call
  \item a flood, an earthquake, a tornado, a drought, to flood
  \item a disaster, an explosion, to explode, to collapse, panic
  \item to choke, to survive, to bandage, to burn oneself, to cut oneself
\end{itemize}

\begin{tcolorbox}[breakable,colback=teal!4,colframe=teal!35!black,title={Before you move on}]
A fire? (\textbf{\ru{пожар}}.) Smoke? (\textbf{\ru{дым}}.) To burn? (\textbf{\ru{гореть}}.) To catch fire? (\textbf{\ru{загораться}}.) A fire extinguisher? (\textbf{\ru{огнетушитель}}.) A rescuer? (\textbf{\ru{спасатель}}.) To save? (\textbf{\ru{спасать}}.) Rescue? (\textbf{\ru{спасение}}.) To evacuate? (\textbf{\ru{эвакуировать}}.) An evacuation? (\textbf{\ru{эвакуация}}.) An injury? (\textbf{\ru{травма}}.) A bruise? (\textbf{\ru{ушиб}}.) Consciousness? (\textbf{\ru{сознание}}.) To be hurt? (\textbf{\ru{пострадать}}.) To wound? (\textbf{\ru{ранить}}.) A hit? (\textbf{\ru{наезд}}.) An injured person? (\textbf{\ru{пострадавший}}.) To crash? (\textbf{\ru{разбиваться}}.) To come out unhurt? (\textbf{\ru{уцелеть}}.) To overturn? (\textbf{\ru{перевернуться}}.) A stretcher? (\textbf{\ru{носилки}}.) A siren? (\textbf{\ru{сирена}}.) An alarm? (\textbf{\ru{тревога}}.) The person on duty? (\textbf{\ru{дежурный}}.) To call? (\textbf{\ru{вызывать}}.) A flood? (\textbf{\ru{наводнение}}.) An earthquake? (\textbf{\ru{землетрясение}}.) A tornado? (\textbf{\ru{смерч}}.) A drought? (\textbf{\ru{засуха}}.) To flood? (\textbf{\ru{затоплять}}.) A disaster? (\textbf{\ru{катастрофа}}.) An explosion? (\textbf{\ru{взрыв}}.) To explode? (\textbf{\ru{взрываться}}.) To collapse? (\textbf{\ru{обрушиться}}.) Panic? (\textbf{\ru{паника}}.) To choke? (\textbf{\ru{задыхаться}}.) To survive? (\textbf{\ru{выживать}}.) To bandage? (\textbf{\ru{перевязывать}}.) To burn oneself? (\textbf{\ru{обжечься}}.) To cut oneself? (\textbf{\ru{порезаться}}.)
\end{tcolorbox}
